\documentclass[11pt,a4paper]{article}
\usepackage[utf8]{inputenc}
\usepackage[T1]{fontenc}
\usepackage[magyar]{babel}
\usepackage{amsmath,amssymb}
\usepackage{geometry}
\usepackage{tcolorbox}
\usepackage{xcolor}
\usepackage{tikz}

\geometry{top=2cm, bottom=2cm, left=2cm, right=2cm}

\title{Hatványozás azonosságai Dolgozat_Dolgozat}
\author{AI LaTeX Studio}
\date{\today}

%%% Megjelenítéshez automatikusan kiegészített csomagok %%%
\usepackage{amsfonts, amsthm}
\usepackage{pgfplots}
\pgfplotsset{compat=1.18}
\usepackage{graphicx}
\usepackage{enumitem}
\usepackage{booktabs}
\usepackage{float}
\usepackage{caption}
\usepackage{microtype}

\usepackage{hyperref}

\begin{document}

% =========================================================================
% "A" CSOPORT DOLGOZAT
% =========================================================================

\noindent
\begin{minipage}[t]{0.45\textwidth}
  \vspace{0pt}
  {\Large\textbf{Hatványozás azonosságai Dolgozat}}\\[2mm]
  {\large\textbf{„A” CSOPORT}}\\[2mm]
  \textbf{Név:} \underline{\hspace{4.5cm}}\\[2mm]
  \textbf{Osztály:} \underline{\hspace{1.5cm}} \quad \textbf{Dátum:} \underline{\hspace{2cm}}
\end{minipage}
\hfill
\begin{minipage}[t]{0.53\textwidth}
  \vspace{0pt}
  \begin{flushright}
  {\small
  \begin{tabular}{|l|c|c|c|c|c||c||c|}
    \hline
    \textbf{Feladat} & 1. & 2. & 3. & 4. & 5. & \textbf{Össz.} & \textbf{Érdemjegy} \\ \hline
    \textbf{Max. pont} & 8 & 8 & 8 & 8 & 8 & 40 & \\ \cline{1-7}
    \textbf{Elért pont} & & & & & & & \\ \hline
  \end{tabular}
  }
  \end{flushright}
\end{minipage}
\vspace{2mm}
\hrule
\vspace{4mm}

\noindent
\textit{Munkaidő: 45 perc. Minden feladatnál követelmény a levezetés, a részeredmények feltüntetése, valamint a végeredmény legegyszerűbb alakban történő megadása. Zsebszámológép használható.}

\vspace{3mm}

\begin{enumerate}
    \item \textbf{Alapvető hatványozási azonosságok alkalmazása} \hfill \textbf{[8 pont]}
    
    Írd fel egyetlen hatvány alakjában és számítsd ki az alábbi kifejezések értékét hatványalap és hatványkitevő azonosságainak felhasználásával:
    \begin{enumerate}
        \item $\dfrac{2^5 \cdot 2^{-3}}{2^2}$ \hfill \textit{[2 pont]}
        \item $(3^2)^3 \cdot 3^{-4}$ \hfill \textit{[2 pont]}
        \item $\dfrac{5^6 \cdot 25^2}{125^2}$ \hfill \textit{[4 pont]}
    \end{enumerate}

    \vspace{2mm}
    \item \textbf{Algebrai kifejezések egyszerűsítése hatványokkal} \hfill \textbf{[8 pont]}
    
    Egyszerűsítsd az alábbi algebrai kifejezéseket, és add meg a feltételeket, amelyek mellett a kifejezés értelmezve van ($x \neq 0, y \neq 0$):
    \begin{enumerate}
        \item $\left(\frac{2x^3}{y^2}\right)^{-2} \cdot \left(\frac{x^{-1}}{4y^3}\right)$ \hfill \textit{[4 pont]}
        \item $\dfrac{(x^2 y^{-3})^3 \cdot x^{-4}}{x^{-2} y^5}$ \hfill \textit{[4 pont]}
    \end{enumerate}

    \vspace{2mm}
    \item \textbf{Normálalak és számítások nagyságrenddel} \hfill \textbf{[8 pont]}
    
    Végezd el a következő műveleteket, és az eredményt add meg normálalakban:
    \begin{enumerate}
        \item $(3,2 \cdot 10^5) \cdot (4,5 \cdot 10^{-2})$ \hfill \textit{[3 pont]}
        \item $\dfrac{8,4 \cdot 10^7}{2,1 \cdot 10^3}$ \hfill \textit{[2 pont]}
        \item $(1,5 \cdot 10^4) + (2,3 \cdot 10^3)$ \hfill \textit{[3 pont]}
    \end{enumerate}

    \vspace{2mm}
    \item \textbf{Exponenciális egyenletek megoldása} \hfill \textbf{[8 pont]}
    
    Oldód meg a következő exponenciális egyenleteket a valós számok halmazán az alapok azonosságára visszavezetve:
    \begin{enumerate}
        \item $2^{3x-1} = 32$ \hfill \textit{[4 pont]}
        \item $9^{x+1} = \dfrac{1}{27^x}$ \hfill \textit{[4 pont]}
    \end{enumerate}

    \vspace{2mm}
    \item \textbf{Gondolkodtató, komplex feladat} \hfill \textbf{[8 pont]}
    
    Adott a következő kifejezés, ahol $n \in \mathbb{N}$:
    $$A_n = \dfrac{2^{n+2} + 2^n}{2^{n+1} - 2^{n-1}}$$
    \begin{enumerate}
        \item Hozd az $A_n$ kifejezést a lehető simplest alakba kiemelés és egyszerűsítés segítségével! \hfill \textit{[5 pont]}
        \item Határozd meg az $A_{2024}$ értékét! \hfill \textit{[3 pont]}
    \end{enumerate}
\end{enumerate}

\vfill
\noindent
\begin{center}
\small
\textbf{Ponthatárok:} \quad 1 (elégtelen): 0--15 p \quad 2 (elégséges): 16--21 p \quad 3 (közepes): 22--27 p \quad 4 (jó): 28--33 p \quad 5 (jeles): 34--40 p
\end{center}


\newpage
% =========================================================================
% "B" CSOPORT DOLGOZAT
% =========================================================================

\noindent
\begin{minipage}[t]{0.45\textwidth}
  \vspace{0pt}
  {\Large\textbf{Hatványozás azonosságai Dolgozat}}\\[2mm]
  {\large\textbf{„B” CSOPORT}}\\[2mm]
  \textbf{Név:} \underline{\hspace{4.5cm}}\\[2mm]
  \textbf{Osztály:} \underline{\hspace{1.5cm}} \quad \textbf{Dátum:} \underline{\hspace{2cm}}
\end{minipage}
\hfill
\begin{minipage}[t]{0.53\textwidth}
  \vspace{0pt}
  \begin{flushright}
  {\small
  \begin{tabular}{|l|c|c|c|c|c||c||c|}
    \hline
    \textbf{Feladat} & 1. & 2. & 3. & 4. & 5. & \textbf{Össz.} & \textbf{Érdemjegy} \\ \hline
    \textbf{Max. pont} & 8 & 8 & 8 & 8 & 8 & 40 & \\ \cline{1-7}
    \textbf{Elért pont} & & & & & & & \\ \hline
  \end{tabular}
  }
  \end{flushright}
\end{minipage}
\vspace{2mm}
\hrule
\vspace{4mm}

\noindent
\textit{Munkaidő: 45 perc. Minden feladatnál követelmény a levezetés, a részeredmények feltüntetése, valamint a végeredmény legegyszerűbb alakban történő megadása. Zsebszámológép használható.}

\vspace{3mm}

\begin{enumerate}
    \item \textbf{Alapvető hatványozási azonosságok alkalmazása} \hfill \textbf{[8 pont]}
    
    Írd fel egyetlen hatvány alakjában és számítsd ki az alábbi kifejezések értékét hatványalap és hatványkitevő azonosságainak felhasználásával:
    \begin{enumerate}
        \item $\dfrac{3^6 \cdot 3^{-2}}{3^3}$ \hfill \textit{[2 pont]}
        \item $(2^3)^2 \cdot 2^{-5}$ \hfill \textit{[2 pont]}
        \item $\dfrac{2^8 \cdot 4^3}{8^3}$ \hfill \textit{[4 pont]}
    \end{enumerate}

    \vspace{2mm}
    \item \textbf{Algebrai kifejezések egyszerűsítése hatványokkal} \hfill \textbf{[8 pont]}
    
    Egyszerűsítsd az alábbi algebrai kifejezéseket, és add meg a feltételeket, amelyek mellett a kifejezés értelmezve van ($a \neq 0, b \neq 0$):
    \begin{enumerate}
        \item $\left(\frac{3a^2}{b^3}\right)^{-2} \cdot \left(\frac{a^{-2}}{3b^2}\right)$ \hfill \textit{[4 pont]}
        \item $\dfrac{(a^3 b^{-2})^2 \cdot a^{-3}}{a^{-1} b^4}$ \hfill \textit{[4 pont]}
    \end{enumerate}

    \vspace{2mm}
    \item \textbf{Normálalak és számítások nagyságrenddel} \hfill \textbf{[8 pont]}
    
    Végezd el a következő műveleteket, és az eredményt add meg normálalakban:
    \begin{enumerate}
        \item $(2,4 \cdot 10^4) \cdot (5,0 \cdot 10^{-3})$ \hfill \textit{[3 pont]}
        \item $\dfrac{9,6 \cdot 10^8}{3,2 \cdot 10^4}$ \hfill \textit{[2 pont]}
        \item $(2,1 \cdot 10^5) + (4,2 \cdot 10^4)$ \hfill \textit{[3 pont]}
    \end{enumerate}

    \vspace{2mm}
    \item \textbf{Exponenciális egyenletek megoldása} \hfill \textbf{[8 pont]}
    
    Oldód meg a következő exponenciális egyenleteket a valós számok halmazán az alapok azonosságára visszavezetve:
    \begin{enumerate}
        \item $3^{2x-1} = 243$ \hfill \textit{[4 pont]}
        \item $4^{x-1} = \dfrac{1}{8^x}$ \hfill \textit{[4 pont]}
    \end{enumerate}

    \vspace{2mm}
    \item \textbf{Gondolkodtató, komplex feladat} \hfill \textbf{[8 pont]}
    
    Adott a következő kifejezés, ahol $n \in \mathbb{N}$:
    $$B_n = \dfrac{3^{n+2} + 3^n}{3^{n+1} - 3^{n-1}}$$
    \begin{enumerate}
        \item Hozd a $B_n$ kifejezést a lehető simplest alakba kiemelés és egyszerűsítés segítségével! \hfill \textit{[5 pont]}
        \item Határozd meg a $B_{2024}$ értékét! \hfill \textit{[3 pont]}
    \end{enumerate}
\end{enumerate}

\vfill
\noindent
\begin{center}
\small
\textbf{Ponthatárok:} \quad 1 (elégtelen): 0--15 p \quad 2 (elégséges): 16--21 p \quad 3 (közepes): 22--27 p \quad 4 (jó): 28--33 p \quad 5 (jeles): 34--40 p
\end{center}

\end{document}